\chapter{taskman --- The Central System Server}

\section{Role}

\texttt{taskman} is the root server process in QSOE.  It is analogous
to QNX Neutrino's \texttt{procnto}: a single user-space process that
combines the responsibilities of a process manager, memory manager,
path manager, and system service provider.

Because seL4 provides only the bare capability and IPC mechanisms,
\texttt{taskman} is responsible for all higher-level OS policy.

\section{Startup}

After sel4runtime hands control to \texttt{main()}, taskman:

\begin{enumerate}
    \item Inspects \texttt{seL4\_BootInfo} to discover available untyped
          memory and initial capabilities.
    \item Creates its main channel (an seL4 Endpoint) and begins
          listening for incoming messages.
    \item Initializes internal subsystems: memory allocator, process
          table, pathname space.
\end{enumerate}

\section{Services}

\texttt{taskman} multiplexes several service domains over a single
channel, dispatching on message type:

\begin{table}[ht]
\centering
\begin{tabular}{@{}lll@{}}
    \toprule
    \textbf{Domain}   & \textbf{Operations}
                       & \textbf{QNX equivalent} \\
    \midrule
    Memory             & \texttt{mmap()}, \texttt{munmap()}, shared memory
                       & \texttt{procnto} memmgr \\
    Process            & \texttt{ProcessCreate()}, \texttt{ProcessTerminate()}
                       & \texttt{procnto} process manager \\
    Pathname           & path registration, resolution, \texttt{open()}
                       & \texttt{procnto} pathmgr \\
    System             & \texttt{sysinfo()}, \texttt{sysshutdown()},
                         \texttt{ioport\_access\_grant()}
                       & various \texttt{procnto} calls \\
    \bottomrule
\end{tabular}
\caption{taskman service domains}
\label{tab:taskman-services}
\end{table}

\section{Message-passing model}

All interaction with \texttt{taskman} follows the QNX
Send/Receive/Reply pattern:

\begin{enumerate}
    \item A client calls \texttt{MsgSend()} (maps to \texttt{seL4\_Call}),
          which blocks the sender.
    \item \texttt{taskman} calls \texttt{MsgReceive()} (maps to
          \texttt{seL4\_Recv}) to dequeue the next message.
    \item \texttt{taskman} processes the request and calls
          \texttt{MsgReply()} (maps to \texttt{seL4\_Reply}) to unblock
          the sender with the result.
\end{enumerate}

Asynchronous notifications (pulses) use \texttt{MsgSendPulse()},
mapped to seL4 Notification objects, for interrupt delivery and
lightweight signaling.
